\section{Introduction}
\label{sec:introduction}

Compactness is perhaps the most widely invoked redistricting criterion: it appears in the constitutions or statutes of 37 states, is regularly cited by courts reviewing redistricting maps, and has been proposed as a proxy for partisan neutrality on the theory that compact districts cannot be precisely gerrymandered.\footnote{See \citet{polsby1991third}; \citet{grofman1985criteria}.} Yet the empirical literature on algorithmic redistricting compactness has focused almost entirely on the 435 congressional districts---a small fraction of the approximately 7,383 state legislative districts drawn in each redistricting cycle.

This focus is understandable. Congressional redistricting is the most politically consequential form of redistricting, and the data infrastructure for congressional maps (TIGER shapefiles, Census Redistricting Data Summary Files) is more uniformly available than state legislative data. But the methodological consequence is that we know much more about how algorithms perform at drawing 4--53 districts per state than we know about how they perform at drawing 40--400 districts per state.

This paper fills that gap. We present a systematic empirical comparison of algorithmic compactness at two scales: congressional (mean 8.7 districts per state) and state house (mean 99.3 districts per state). The comparison exploits a simple mathematical relationship: for a fixed-area state, mean district perimeter scales as $O(1/\sqrt{k})$ as the number of districts $k$ increases, which implies that Polsby-Popper ratios---which involve perimeter squared in the denominator---improve as the number of districts increases. This relationship is a mathematical consequence of the geometry of partitioning a bounded area, not a property of any particular algorithm.

We confirm this theoretical prediction empirically. Across 50 states and three census years (2000, 2010, 2020), algorithmic state house maps (drawn at tract resolution) achieve a mean Polsby-Popper ratio of 0.388 versus 0.361 for congressional maps---a 7.5\% improvement. At block group resolution, which traces district boundaries more precisely, the state house advantage widens to 11\%: mean PP of 0.401 versus 0.361.

We also examine the role of geographic resolution---tract versus block group---in determining compactness. Block groups are smaller geographic units (mean 1,500 persons) than census tracts (mean 4,000 persons), allowing the algorithm to trace district boundaries more precisely, following natural geographic features more closely and producing districts with shorter perimeters relative to their area. The resolution effect averages 0.020 PP units across all states, but varies substantially, with the largest resolution gains in geographically complex states (Alaska, Hawaii, coastal states with irregular boundaries).

Finally, we compare algorithmic state house maps to enacted state legislative maps using GIS data for 35 states where enacted maps are available. Algorithmic maps outperform enacted maps by a mean of 0.062 PP units (17\%), with the largest gains in states where enacted maps have been identified as gerrymandered by courts or independent analysts. This comparison establishes that the compactness advantage of algorithmic redistricting is not merely a feature of the algorithm's performance against a random baseline---it represents a genuine improvement over the maps that human mapmakers actually draw.

The paper is organized as follows. Section~\ref{sec:math} develops the mathematical intuition for why smaller districts should be more compact. Section~\ref{sec:empirical} presents the empirical comparison of state house and congressional PP ratios across all 50 states. Section~\ref{sec:resolution} examines the resolution effect. Section~\ref{sec:enacted} compares algorithmic state house maps to enacted maps. Section~\ref{sec:conclusion} concludes.

\subsection{Significance for State Legislative Redistricting Reform}

The state legislative compactness finding has practical significance beyond academic interest. State legislative redistricting is both more numerous and more consequential for daily governance than congressional redistricting: state legislatures draw their own congressional maps, control state-level policy, and—through the Elections Clause---govern the rules of federal elections. A finding that algorithmic maps achieve substantially better compactness at state legislative scale than enacted maps do strengthens the case for algorithmic adoption at the state level.

Moreover, many states impose stronger compactness requirements for state legislative districts than for congressional districts. Iowa's redistricting statute, for example, explicitly ranks compactness as a primary criterion for state legislative maps. The finding that the algorithm performs better on compactness at state legislative scale makes it a stronger candidate for adoption in compactness-stringent states.
